\documentclass[preview]{standalone}

\usepackage{amsmath}
\usepackage{amssymb}
\usepackage{stellar}
\usepackage{definitions}

\begin{document}

\id{psicologia-memoria-lungo-termine}
\genpage

\section{Memoria a lungo termine}

\begin{snippetdefinition}{memoria-lungo-termine-definition}{Memoria a lungo termine}
    La \emph{memoria a lungo termine}, o \emph{MLT}, è il magazzino delle
    esperienze, degli eventi, delle informazioni e delle capacità che una
    persona possiede su di sé e sul mondo, acquisite attraverso la memoria
    sensoriale e la memoria di lavoro.
\end{snippetdefinition}

\includesnpt[width=58\%,src=/snippet/static-psychology/memory-systems-model.jpg]{centered-img}

\begin{snippetdefinition}{indizio-recupero-definition}{Indizio di recupero}
    Un \emph{indizio di recupero} è un elemento contestuale, esterno o interno,
    che permette di recuperare ricordi immagazzinati nella memoria a lungo
    termine.
\end{snippetdefinition}

\begin{snippetdefinition}{richiamo-definition}{Richiamo}
    Il \emph{richiamo} è una tecnica di recupero della memoria in cui il
    soggetto deve riprodurre un'informazione a cui è stato precedentemente
    esposto.
\end{snippetdefinition}

\begin{snippetdefinition}{riconoscimento-definition}{Riconoscimento}
    Il \emph{riconoscimento} è la realizzazione che un dato evento o stimolo è
    già stato esperito in passato.
\end{snippetdefinition}

\begin{snippetexample}{testimonianza-recupero-riconoscimento-example}{Testimonianza}
    Se la polizia chiede a una vittima di ricordare alcune caratteristiche
    distintive dell'aggressore, usa un metodo basato sul richiamo. Se invece
    mostra alla vittima fotografie di sospetti, usa un metodo basato sul
    riconoscimento.
\end{snippetexample}

\begin{snippetdefinition}{memoria-episodica-definition}{Memoria episodica}
    La \emph{memoria episodica} conserva eventi specifici di cui si fa
    esperienza personalmente. Per riattivare questi ricordi sono spesso
    importanti indizi di contesto legati al momento in cui l'evento è accaduto.
\end{snippetdefinition}

\begin{snippetdefinition}{memoria-semantica-definition}{Memoria semantica}
    La \emph{memoria semantica} è una memoria generica e categoriale, relativa
    al significato delle parole, dei concetti, dei simboli e alle relazioni tra
    essi.
\end{snippetdefinition}

\begin{snippetexample}{memoria-episodica-semantica-example}{Memoria episodica e semantica}
    Il ricordo del primo bacio è un esempio di memoria episodica, perché è
    legato al contesto personale in cui l'evento è avvenuto. Conoscenze come la
    capitale della Francia o una costante fisica appartengono invece alla
    memoria semantica.
\end{snippetexample}

\begin{snippetdefinition}{specificita-contesto-codifica-definition}{Specificità del contesto di codifica}
    Il principio della \emph{specificità del contesto di codifica} afferma che i
    ricordi si attivano più rapidamente quando il contesto di recupero è coerente
    con il contesto in cui sono stati codificati.
\end{snippetdefinition}

\begin{snippetdefinition}{memoria-contesto-dipendente-definition}{Memoria contesto-dipendente}
    La \emph{memoria contesto-dipendente} è il fenomeno per cui il recupero di
    un'informazione è facilitato quando il contesto in cui viene recuperata è
    simile a quello in cui è stata codificata.
\end{snippetdefinition}

\begin{snippetexample}{riconoscimento-fuori-contesto-example}{Riconoscimento fuori contesto}
    Può capitare di vedere in una stanza affollata una persona che si sa di aver
    già incontrato, senza riuscire a ricordare dove. La difficoltà nasce dal
    fatto che la persona non era mai stata vista in quel particolare contesto.
\end{snippetexample}

\begin{snippetdefinition}{effetto-posizione-seriale-definition}{Effetto di posizione seriale}
    L'\emph{effetto di posizione seriale} è il fenomeno per cui, in una lista,
    vengono ricordati meglio i primi elementi e gli ultimi elementi, mentre gli
    elementi intermedi sono ricordati con maggiore difficoltà.
\end{snippetdefinition}

\begin{snippet}{priorita-recenza}
    L'effetto di posizione seriale comprende l'\emph{effetto priorità}, cioè il
    miglior ricordo dei primi item, e l'\emph{effetto recenza}, cioè il miglior
    ricordo degli ultimi item.
\end{snippet}

\includesnpt[width=62\%,src=/snippet/static-psychology/serial-position-curve.jpg]{centered-img}

\begin{snippetdefinition}{differenziazione-contestuale-definition}{Differenziazione contestuale}
    La \emph{differenziazione contestuale} è la diversità del contesto in cui le
    informazioni sono state apprese, che può renderle più o meno distinguibili e
    quindi più o meno facili da recuperare.
\end{snippetdefinition}

\includesnpt[width=58\%,src=/snippet/static-psychology/contextual-differentiation.jpg]{centered-img}

\begin{snippetdefinition}{teoria-livelli-elaborazione-definition}{Teoria dei livelli di elaborazione}
    La \emph{teoria dei livelli di elaborazione} sostiene che quanto più è
    profondo il livello a cui un'informazione viene elaborata, tanto più è
    probabile che essa venga recuperata dalla memoria.
\end{snippetdefinition}

\begin{snippetexample}{livelli-elaborazione-example}{Livelli di elaborazione}
    Consideriamo la parola \texttt{CANE}. Si può osservare che è scritta in
    maiuscolo, che termina con \texttt{-ne}, oppure che indica un animale a
    quattro zampe. Ognuna di queste osservazioni richiede un livello di
    elaborazione via via più profondo.
\end{snippetexample}

\begin{snippetdefinition}{priming-definition}{Priming}
    Il \emph{priming} è un fenomeno per cui uno stimolo precedente influenza la
    risposta a uno stimolo successivo, anche in assenza di una correlazione
    evidente tra i due.
\end{snippetdefinition}

\end{document}
